\documentclass[11pt,letterpaper]{article}
\usepackage{physical_ai_legislation}
\usepackage[backend=biber,style=numeric,sorting=none]{biblatex}
\addbibresource{physical_ai_building_code.bib}

\title{%
  \textbf{Physical AI Oncology Clinical Trial Facility}\\
  \textbf{Building Code Standards}\\[6pt]
  \large Construction, Structural, Mechanical, Electrical,\\
  and Technology Infrastructure Requirements\\[6pt]
  \normalsize Supplemental standards to California Building Standards Code\\
  (Title 24) for Physical AI oncology clinical trial site construction.
}
\author{%
  Kevin Kawchak\\
  CEO, ChemicalQDevice\\
  \texttt{kevink@chemicalqdevice.com}
}
\date{March 24, 2026}

\begin{document}
\maketitle
\thispagestyle{fancy}

\begin{center}
\footnotesize
\textit{This draft is independent and is not endorsed, sponsored, or approved
by any trial sponsor, CRO, site, IRB, regulator, or medical society;
and was adapted using Claude Code Opus 4.6.}
\end{center}

\vspace{1em}
\tableofcontents
\newpage

\section{General Requirements}

\subsection{Scope}

These building code standards supplement the California Building Standards Code
(Title 24, Parts 1 through 12) with additional requirements specific to Physical
AI oncology clinical trial facilities. These standards apply to new construction
of facilities designed to house 10 categories of Physical AI robotic systems
operating autonomously on a 24-hour, 7-day-per-week schedule serving up to 180
oncology patients per day~\cite{sitespec2026}.

\subsection{Facility Size and Configuration}

\begin{enumerate}[label=(\alph*)]
\item The minimum gross floor area shall be 85,000 square feet, with a maximum
  of 100,000 square feet, distributed across the clinical wings, support areas,
  and infrastructure spaces specified in these standards~\cite{sitespec2026}.

\item The facility shall accommodate nine clinical wings corresponding to the
  following Physical AI system deployment zones:
  \begin{enumerate}[label=(\arabic*)]
  \item Three surgical suites for surgical robots (SURG-01, SURG-02, SURG-03).
  \item Four cobot stations for collaborative biopsy robots.
  \item Three radiotherapy vaults for RT positioning robots.
  \item Two needle-placement suites for robotic needle-placement systems.
  \item Five social companion stations for pediatric and adult companion
    interactions.
  \item Three humanoid stations for rehabilitation engagement.
  \item Three RT motion-tracking vaults (co-located with RT positioning).
  \item Four imaging bays for imaging assistant robots.
  \item Two steerable needle suites.
  \item One rehabilitation exoskeleton bay with two active stations and one
    reserved station.
  \end{enumerate}

\item Support areas shall include a robot-accessible pharmacy, server room,
  maintenance bay, and decontamination area~\cite{newtrial2026}.
\end{enumerate}

\section{Structural Requirements}

\subsection{Floor Load Capacity}

\begin{enumerate}[label=(\alph*)]
\item Surgical suites shall be designed for a minimum floor live load of 150
  pounds per square foot to support surgical robot systems (including the
  da Vinci Xi surgical system at approximately 1,800 pounds), instrument carts,
  patient tables, and support equipment~\cite{patientinstructions2026}.

\item Radiotherapy vaults shall be designed for a minimum floor live load of
  200 pounds per square foot to support linear accelerator systems,
  radiotherapy positioning robots, and radiation shielding components.

\item Imaging bays shall be designed for a minimum floor live load of 125
  pounds per square foot to support imaging equipment, robotic ultrasound
  systems, and CT scanners used for needle-placement guidance.

\item Rehabilitation exoskeleton areas shall include floor-mounted parallel
  bars with anchorage rated for 500 pounds lateral force, supporting patients
  during exoskeleton-assisted gait
  training~\cite{patientinstructions2026}.

\item Server rooms shall be designed for a minimum floor live load of 250
  pounds per square foot to support server rack installations for AI
  computation infrastructure.

\item General clinical circulation areas shall maintain a minimum floor live
  load of 100 pounds per square foot, accommodating hospital beds,
  wheelchairs, and mobile robot systems simultaneously.
\end{enumerate}

\subsection{Vibration Control}

\begin{enumerate}[label=(\alph*)]
\item Surgical suites and needle-placement suites shall meet VC-E vibration
  criteria (peak velocity 3.1 micrometers per second) to support the
  submillimeter precision required by surgical robots and needle-placement
  systems~\cite{newtrial2026}.

\item Imaging bays shall meet VC-D vibration criteria (peak velocity 6.3
  micrometers per second) to support diagnostic imaging quality.

\item Radiotherapy vaults shall meet VC-C vibration criteria (peak velocity
  12.5 micrometers per second) to support accurate beam alignment.

\item Vibration isolation joints shall separate clinical wings from building
  mechanical systems, loading areas, and parking structures.
\end{enumerate}

\section{Mechanical Systems (HVAC)}

\subsection{Surgical Suite Air Handling}

\begin{enumerate}[label=(\alph*)]
\item Surgical suites shall maintain ISO Class 7 (10,000 particles per cubic
  foot) or better air cleanliness using HEPA filtration with a minimum of 20
  air changes per hour~\cite{sitespec2026}.

\item Surgical suites shall maintain positive pressure of at least 2.5 Pascals
  relative to adjacent spaces.

\item Temperature control shall maintain 68 to 73 degrees Fahrenheit with
  humidity between 30 and 60 percent relative humidity.

\item Independent air handling units shall serve each surgical suite to prevent
  cross-contamination and allow maintenance without affecting other suites.
\end{enumerate}

\subsection{Radiotherapy Vault Ventilation}

\begin{enumerate}[label=(\alph*)]
\item Radiotherapy vaults shall have dedicated exhaust systems capable of
  removing ozone generated by high-energy radiation beams.

\item A minimum of 6 air changes per hour shall be maintained during treatment
  operations, with 10 air changes per hour during purge cycles between patients.

\item Exhaust air from radiotherapy vaults shall not be recirculated to other
  building areas.
\end{enumerate}

\subsection{Server Room Cooling}

\begin{enumerate}[label=(\alph*)]
\item Server rooms supporting AI computation infrastructure shall maintain
  temperature between 64 and 75 degrees Fahrenheit with redundant cooling
  systems providing N+1 capacity~\cite{sitespec2026}.

\item Cooling systems shall reject a minimum heat load of 200 kilowatts,
  calculated from the server room's portion of the facility's total 800 to
  1,200 kilowatt electrical load.

\item Server room cooling shall be on a dedicated HVAC zone independent of
  clinical space climate control.

\item Hot aisle and cold aisle containment shall be used to maximize cooling
  efficiency.
\end{enumerate}

\section{Electrical Systems}

\subsection{Power Supply}

\begin{enumerate}[label=(\alph*)]
\item The facility shall be served by dual utility power feeds from separate
  substations or feeders, providing full redundancy for the 800 to 1,200
  kilowatt total electrical load~\cite{sitespec2026}.

\item Automatic transfer switches shall provide seamless transition between
  utility feeds, with a maximum transfer time of 10 milliseconds for critical
  clinical systems.

\item Emergency generator capacity shall be sufficient to power all critical
  clinical systems (surgical suites, radiotherapy vaults, server rooms,
  life safety systems) for a minimum of 72 hours at full load.

\item Generator fuel storage shall accommodate 96 hours of operation with
  automatic fuel management and monitoring.
\end{enumerate}

\subsection{Uninterruptible Power Supply (UPS)}

\begin{enumerate}[label=(\alph*)]
\item All surgical suites shall be served by dedicated UPS systems providing a
  minimum of 30 minutes of battery backup at full surgical robot operating
  load, sufficient to safely complete or abort any procedure in progress.

\item Server rooms shall be served by UPS systems providing a minimum of 15
  minutes of battery backup, sufficient for orderly shutdown of AI computation
  systems and data preservation.

\item All Physical AI system charging and standby stations shall be served by
  UPS-backed circuits to maintain robot readiness during power transitions.

\item UPS systems shall be tested monthly under load, with results recorded
  in the facility maintenance log.
\end{enumerate}

\subsection{Electrical Distribution}

\begin{enumerate}[label=(\alph*)]
\item Dedicated electrical panels shall serve each clinical wing, with
  sufficient circuit capacity for all Physical AI system types deployed in
  that wing plus 25 percent spare capacity.

\item Isolated ground circuits shall serve all Physical AI systems and medical
  imaging equipment to prevent electromagnetic interference.

\item Power quality monitoring systems shall continuously track voltage, current,
  frequency, harmonics, and transients on all circuits serving Physical AI
  systems, with automated alerting for out-of-specification conditions.
\end{enumerate}

\section{Radiation Shielding}

\begin{enumerate}[label=(\alph*)]
\item Radiotherapy vault shielding shall comply with NCRP Report 151 and
  California Radiologic Health Branch requirements, designed for the maximum
  beam energy and workload planned for the facility.

\item Shielding design shall account for the Physical AI positioning robot's
  range of motion, ensuring that all achievable patient positions result in
  adequate shielding of adjacent spaces.

\item The primary barrier shall be designed for the full workload of the
  radiotherapy system, and secondary barriers shall account for leakage and
  scatter radiation.

\item Maze entrances shall be designed to reduce radiation levels at the maze
  entrance to below 0.02 milliSieverts per hour during beam-on conditions.

\item Shielding calculations shall be reviewed and certified by a qualified
  medical physicist before construction begins.
\end{enumerate}

\section{Technology Infrastructure}

\subsection{Network Architecture}

\begin{enumerate}[label=(\alph*)]
\item The facility shall provide a minimum of 10 gigabits per second internal
  network bandwidth using redundant fiber-optic backbone
  infrastructure~\cite{sitespec2026}.

\item Clinical networks, administrative networks, and Physical AI system control
  networks shall be physically or logically segmented with firewall separation
  between segments.

\item Each clinical wing shall have dedicated network switches with redundant
  uplinks to the core network.

\item Wireless coverage (Wi-Fi 6E or later) shall be provided throughout all
  patient-accessible areas for patient portal access and companion robot
  connectivity.
\end{enumerate}

\subsection{Robot Charging and Docking Infrastructure}

\begin{enumerate}[label=(\alph*)]
\item Each clinical wing shall include dedicated robot charging and docking
  stations rated for the specific power requirements of the Physical AI
  systems deployed in that wing.

\item Charging stations shall include automated connection systems to minimize
  human intervention during shift changes and standby periods.

\item Floor-mounted guide paths or ceiling-mounted positioning references shall
  facilitate autonomous robot navigation between clinical stations, charging
  stations, and maintenance areas.

\item The facility layout shall support the patient flow demonstrated in the
  simulation, where peak utilization reached 21 of 29 robots active
  simultaneously during hour 09, with an average of 8.1 of 29 active across
  24 hours~\cite{newtrial2026}.
\end{enumerate}

\section{Life Safety and Egress}

\begin{enumerate}[label=(\alph*)]
\item The facility shall comply with California Fire Code (Title 24, Part 9)
  and NFPA 101 Life Safety Code requirements for health care occupancies.

\item Emergency egress paths shall be designed to accommodate patients on
  hospital beds and in wheelchairs, with corridor widths of at least 8 feet
  in primary clinical circulation areas.

\item Physical AI systems shall be designed to move to designated safe positions
  during fire alarm activation, clearing egress paths and de-energizing
  non-essential systems.

\item Fire suppression systems in surgical suites shall use clean agent (FM-200
  or Novec 1230) to protect sensitive robotic equipment while maintaining
  patient safety.

\item Fire suppression in server rooms shall use clean agent systems with
  pre-action activation to prevent accidental discharge damage to computing
  equipment.
\end{enumerate}

\printbibliography[title={References}]

\end{document}
